\documentclass[11pt]{article}
\usepackage[margin=1in]{geometry}
\usepackage{amsmath,amssymb,amsthm,mathtools}
\IfFileExists{lmodern.sty}{\usepackage{lmodern}}{}
\usepackage[expansion=false]{microtype}
\usepackage{booktabs,array}
\usepackage{enumitem}
\usepackage{xcolor}
\usepackage[colorlinks=true,linkcolor=blue!50!black,citecolor=green!35!black,urlcolor=blue!50!black]{hyperref}
\usepackage[nameinlink]{cleveref}
\newcolumntype{R}[1]{>{\raggedright\arraybackslash}p{#1}}

\newtheorem{theorem}{Theorem}[section]
\newtheorem{proposition}[theorem]{Proposition}
\newtheorem{lemma}[theorem]{Lemma}
\newtheorem{corollary}[theorem]{Corollary}
\newtheorem{conjecture}[theorem]{Conjecture}
\theoremstyle{definition}
\newtheorem{definition}[theorem]{Definition}
\newtheorem{example}[theorem]{Example}
\newtheorem{question}[theorem]{Question}
\theoremstyle{remark}
\newtheorem{remark}[theorem]{Remark}

\newcommand{\D}{\mathsf{D}}
\newcommand{\Cpl}{\operatorname{Cpl}}
\newcommand{\CP}{\mathbb{CP}}
\newcommand{\TV}{\mathrm{TV}}
\newcommand{\ch}{\operatorname{ch}}
\newcommand{\Coh}{\operatorname{Coh}}

\title{Coupling Polytopes and Toric Geometry:\\
the $2\times3$ Family, the Square Case, Higher Levels and Stability\\[4pt]
\large A companion to \emph{Conditionals, Couplings and Toric Geometry}}
\author{Vinay K. Awasthi\\
Department of Applied Mathematics and Statistics\\
Johns Hopkins University\thanks{Verification script: \texttt{check\_polytope\_companion.py}
(26 checks; exact rational arithmetic except where marked). This version replaces the draft
``Coupling polytopes and the $2\times3$ family''; \cref{app:changes} lists what was corrected and
what was withdrawn.}}
\date{October 5, 2026}

\begin{document}
\maketitle

\begin{abstract}
A coupling of two finite distributions is a point of a transportation polytope. For nondegenerate
marginals this polytope is Delzant \cite{Conditionals}, so it is the moment polytope of a compact
toric symplectic manifold. This companion works the correspondence out in the cases that can be
computed completely, and records how far it goes.
\begin{itemize}[leftmargin=*]
\item \emph{The $2\times3$ family.} The six walls $f_i=g_j$ cut the space of marginals into $18$
  chambers, which fall into $5$ orbits under the symmetry group $S_2\times S_3$. The surfaces that occur
  are $\CP^2$, $\CP^1\times\CP^1$, the Hirzebruch surface $F_1$, $\CP^2$ blown up at two points, and
  the del Pezzo surface of degree $6$. Every facet normal lies in the fan of the last, so every surface
  in the family is the del Pezzo surface or a toric blow-down of it. Crossing a single wall is exactly
  one toric blow-up or blow-down.
\item \emph{The Kähler class.} It is the product coupling read as divisor coefficients,
  $[\omega_{f,g}]=\sum_cf_ig_j[D_c]$. Edge lengths of the coupling polygon are degrees of $\omega$ on
  toric curves, and its area is the symplectic volume.
\item \emph{The square $2\times2$ case.} With the unital whiskering, explicit whiskering defects are
  nonzero, and every defect is a tangent vector of a coupling polytope.
\item \emph{Higher levels.} The Delzant theorem applies at levels two and three. Flattening a level-two
  coupling is affine and commutes with stochastic change of base, but it is neither surjective nor
  injective in general.
\item \emph{Disc potentials.} They are sums over facets, which are cells, not cycles. For the degree-6
  del Pezzo surface the potential has six critical points, all over the product coupling, with critical
  values $6,-2,-2,-2,-3,-3$. Square nondegenerate polytopes with $n\ge3$ give manifolds that are not
  Fano. Optimal transport minimizes a moment-map component, so its optima are torus-fixed points,
  that is, spanning-tree couplings.
\item \emph{Stability.} Each $\omega_{f,g}$ gives a Bridgeland stability condition on the derived category
  of the corresponding surface, with explicit central charges. We state what this does and does not
  connect.
\end{itemize}
\end{abstract}

\tableofcontents

%=====================================================================
\section{Introduction}\label{sec:intro}
%=====================================================================

The companion paper \cite{Conditionals} proves that the polytope $\Cpl(f,g)$ of couplings of two
nondegenerate marginals is Delzant. By Delzant's theorem \cite{Delzant1988} there is therefore a
compact toric symplectic manifold $X_{f,g}$ whose moment map has image $\Cpl(f,g)$. Each interior
coupling is the image of a Lagrangian torus, and each vertex is the image of a fixed point. This note
asks what the correspondence gives in concrete cases.

\Cref{sec:setup} fixes conventions and proves a general formula for the Kähler class.
\Cref{sec:2x3} treats the $2\times3$ family completely. \Cref{sec:square} treats the square $2\times2$
case, the base case for the whiskering questions of \cite[\S7]{Conditionals}. \Cref{sec:levels}
extends the polytope results to second- and third-order states. \Cref{sec:potential} computes disc
potentials and relates optimal transport to the moment map. \Cref{sec:stability} constructs stability
conditions from marginals and states the remaining questions. \Cref{sec:ainfty} records the status of
the $A_\infty$ question. Every numbered claim is checked in \texttt{check\_polytope\_companion.py}.

%=====================================================================
\section{Setup}\label{sec:setup}
%=====================================================================

Let $f\in\D(Y)$ and $g\in\D(Y')$ have full support, with $\lvert Y\rvert=n$ and $\lvert Y'\rvert=m$, and
let
\[
\Cpl(f,g)=\Big\{\alpha\in\mathbb{R}_{\ge0}^{Y\times Y'}:\ \textstyle\sum_{y'}\alpha(y,y')=f(y),\
\sum_y\alpha(y,y')=g(y')\Big\}.
\]
The tangent space is the space of circulations, that is matrices with zero row and column sums. The
integer circulations form a lattice $\Lambda$ of rank $(n-1)(m-1)$ with basis
$e_{ij}=E_{ij}-E_{im}-E_{nj}+E_{nm}$ ($i<n$, $j<m$). Coordinates in this basis are the entries in the
top-left $(n-1)\times(m-1)$ block. The pair $(f,g)$ is \emph{degenerate} if $\sum_If=\sum_Jg$ for some
nonempty proper $I\subsetneq Y$ and $J\subsetneq Y'$. These hyperplanes are the \emph{walls}, and the
components of their complement are the \emph{chambers}.

\begin{theorem}[{\cite[Thm.~4.4]{Conditionals}}]\label{thm:delzant}
If $(f,g)$ is nondegenerate, every vertex of $\Cpl(f,g)$ is supported on a spanning tree of $K_{n,m}$,
the edges at a vertex are the fundamental cycles of that tree, and $\Cpl(f,g)$ is Delzant with respect to
$\Lambda$.
\end{theorem}

\paragraph{Facets and normals.} The facets are among the sets $\{\alpha_c=0\}$, one for each cell
$c=(i,j)$. A cell gives no facet when $f_i+g_j>1$, since then $\alpha_c\ge f_i+g_j-1>0$. The inward normal of
$\{\alpha_c=0\}$ is the coordinate functional $\alpha\mapsto\alpha_c$ restricted to $\Lambda$. In the dual basis
$\{e^*_{ij}\}$ it is
\[
v_{(i,j)}=e^*_{ij},\quad v_{(i,m)}=-\textstyle\sum_je^*_{ij},\quad v_{(n,j)}=-\sum_ie^*_{ij},\quad
v_{(n,m)}=\sum_{i,j}e^*_{ij}\qquad(i<n,\ j<m).
\]
Each $v_c$ is primitive, because some integer circulation has entry $1$ at $c$.

\begin{proposition}[The Kähler class is the product coupling]\label{prop:kahler}
For nondegenerate $(f,g)$, let $D_c$ be the toric divisor of the facet $\{\alpha_c=0\}$. Up to the
normalization of $\omega$, the symplectic class of $X_{f,g}$ is
\[
[\omega_{f,g}]=\sum_{c\ \text{facet}}f_ig_j\,[D_c],\qquad c=(i,j),
\]
and $\omega_{f,g}\cdot C$ is the lattice length of the edge of $\Cpl(f,g)$ that corresponds to a
torus-invariant curve $C$.
\end{proposition}
\begin{proof}
For a Delzant polytope $\{u:\langle u,v_c\rangle+\lambda_c\ge0\}$ the symplectic class is
$\sum_c\lambda_c[D_c]$ \cite{Guillemin1994}. Here the facet functions are
$\alpha_c(u)=\langle u-u_0,v_c\rangle+\alpha_c(u_0)$ for any interior $u_0$, so $\lambda_c=\alpha_c(u_0)$.
Taking $u_0=f\otimes g$ gives the formula. Another base point changes the coefficients by
$\langle w,v_c\rangle$, which is a linear relation among the $[D_c]$. The degree of $\omega$ on a torus-invariant
curve is the lattice length of its edge.
\end{proof}

%=====================================================================
\section{The $2\times3$ family}\label{sec:2x3}
%=====================================================================

Write $f=(f_1,1-f_1)$, $g=(g_1,g_2,1-g_1-g_2)$, $x=\alpha_{11}$ and $y=\alpha_{12}$. The polygon
$\Cpl(f,g)$ is cut out by
\[
x\ge0,\quad y\ge0,\quad x\le g_1,\quad y\le g_2,\quad x+y\le f_1,\quad x+y\ge f_1+g_1+g_2-1,
\]
with inward normals $e_1$, $e_2$, $-e_1$, $-e_2$, $-(e_1+e_2)$ and $e_1+e_2$. The only proper subsets of a
two-point and a three-point set that can have equal sums give the six walls $f_i=g_j$.

\begin{theorem}[Chambers and surfaces]\label{thm:chambers}
The six walls cut the prism $\{0<f_1<1,\ g_1,g_2>0,\ g_1+g_2<1\}$ into $18$ chambers. On each chamber
the polygon type is constant and Delzant, and the surfaces are:
\begin{center}\small
\begin{tabular}{@{}lccll@{}}
\toprule
Polygon & Chambers & Self-intersections $D_c^2$ & Surface & Orbit under $S_2\times S_3$\\
\midrule
triangle & 2 & $(1,1,1)$ & $\CP^2$ & one orbit of size 2\\
quadrilateral & 3 & $(0,0,0,0)$ & $\CP^1\times\CP^1$ & one orbit of size 3\\
quadrilateral & 6 & $(1,0,-1,0)$ & $F_1=\mathrm{Bl}_1\CP^2$ & one orbit of size 6\\
pentagon & 6 & $(0,-1,-1,-1,0)$ & $\mathrm{Bl}_2\CP^2$ & one orbit of size 6\\
hexagon & 1 & $(-1)^6$ & del Pezzo of degree 6 & fixed\\
\bottomrule
\end{tabular}
\end{center}
Here $S_2$ swaps the two rows and $S_3$ permutes the columns.
\end{theorem}
\begin{proof}
For an arrangement of hyperplanes restricted to a convex region, each realizable sign vector is exactly
one convex chamber. A rational grid and random sampling realize the same $18$ sign vectors. In each
chamber the vertices are computed exactly, and the number of vertices and the self-intersections are
constant on all samples. Self-intersections come from the relations $v_{i-1}+v_{i+1}=-D_i^2v_i$
between cyclically ordered normals. Consistency is confirmed by $\sum_iD_i^2=12-3k$ for $k$ divisors.
Delzant means that consecutive normals have determinant $\pm1$. The orbits are computed by letting
$S_2\times S_3$ act on sign vectors.
\end{proof}

\begin{theorem}[Every surface is a blow-down of the del Pezzo surface]\label{thm:blowdown}
For every nondegenerate pair in the $2\times3$ family, $X_{f,g}$ is the del Pezzo surface of degree $6$ or
is obtained from it by toric blow-downs.
\end{theorem}
\begin{proof}
The fan $\Sigma$ of $X_{f,g}$ is smooth and complete. Its rays are among the six rays
$\pm e_1$, $\pm e_2$, $\pm(e_1+e_2)$ of the del Pezzo fan $\Sigma_6$. In cyclic order each ray of $\Sigma_6$ is the
sum of its neighbours. Two of the six rays form a lattice basis exactly when they are at
cyclic distance one or two; rays at distance three are opposite. So adjacent rays $u,w$ of $\Sigma$ are at
distance one or two, and in the second case the omitted ray between them is $u+w$. Inserting $u+w$ is
the star subdivision that blows up the corresponding fixed point \cite{Fulton1993}. Inserting all
omitted rays turns $\Sigma$ into $\Sigma_6$.
\end{proof}

\begin{proposition}[Wall-crossing]\label{prop:wall}
On the wall $f_i=g_j$ three facet lines meet at a point. For example, on $f_1=g_1$ the lines $x=g_1$,
$y=0$ and $x+y=f_1$ meet at $(g_1,0)$. Crossing one wall cuts or restores one corner, which is one toric
blow-up or blow-down. There are no flips, which do not occur for surfaces, and the type does not change
inside a chamber.
\end{proposition}
\begin{proof}
The concurrency is direct. For all $30$ pairs of chambers whose sign vectors differ in one wall, the
vertex counts differ by exactly one. Constancy inside a chamber is \cite[Prop.~4.6]{Conditionals}.
\end{proof}

\begin{example}\label{ex:2x3}
\leavevmode
\begin{enumerate}[label=(\alph*),leftmargin=*]
\item At $(f_1,g_1,g_2)=(0.1,0.5,0.2)$ the polygon is the triangle with vertices $(0,0)$, $(0.1,0)$ and
  $(0,0.1)$, and the surface is $\CP^2$.
\item At $(0.4,0.5,0.3)$ the polygon is the pentagon with vertices $(0,0.2)$, $(0,0.3)$, $(0.1,0.3)$,
  $(0.4,0)$ and $(0.2,0)$, and the surface is $\mathrm{Bl}_2\CP^2$. The straight path to the uniform point
  crosses only the wall $f_1=g_1$; the corner at $(f_1,0)$ is cut and the hexagon appears.
\item At the uniform point $(\tfrac12,\tfrac13,\tfrac13)$ the hexagon has vertices $(0,\tfrac16)$, $(0,\tfrac13)$,
  $(\tfrac16,\tfrac13)$, $(\tfrac13,\tfrac16)$, $(\tfrac13,0)$ and $(\tfrac16,0)$, all six edges have lattice length
  $\tfrac16$, its lattice area is $\tfrac1{12}$, and $[\omega]=-\tfrac16K_X$. This is the unique monotone point of
  the family \cite[Prop.~4.7]{Conditionals}: monotonicity holds at a point, not on a chamber.
\end{enumerate}
\end{example}

\begin{proposition}[Kähler data in the family]\label{prop:kahler2x3}
In all $18$ chambers, $\omega_{f,g}\cdot D_c$ equals the lattice length of the edge $\{\alpha_c=0\}$ and is
positive, so $\omega_{f,g}$ is ample, and $\tfrac12\omega_{f,g}^2$ equals the lattice area of $\Cpl(f,g)$.
\end{proposition}
\begin{proof}
These are \cref{prop:kahler} and the Duistermaat--Heckman formula for the volume. The script verifies
both in exact arithmetic, using the intersection matrix of each fan.
\end{proof}

%=====================================================================
\section{The square $2\times2$ case}\label{sec:square}
%=====================================================================

Let $B=B'=\{0,1\}$. A 1-cell is a stochastic matrix $k(a,b)$ with rows $(1-a,a)$ and $(1-b,b)$, and
Kleisli composition is matrix multiplication. The deterministic 1-cells are the constants $e_0=k(0,0)$
and $e_1=k(1,1)$, the identity $\sigma=k(0,1)$ and the swap $\tau=k(1,0)$. A coupling 2-cell
$\alpha\colon f\Rightarrow g$ is a coupling $\alpha(x)$ of $f(x)$ and $g(x)$ at each $x$. Its cost is
$c(\alpha)=\max_x\Pr_{\alpha(x)}(y\ne y')$. Vertical composition is the gluing lemma, strictly associative
with diagonal identities \cite[Thm.~2.3]{Conditionals}.

\begin{proposition}[Coupling segments and costs]\label{prop:segment}
For $f(x)=(p,1-p)$ and $g(x)=(q,1-q)$, the couplings are
$\left(\begin{smallmatrix}a&p-a\\q-a&1-p-q+a\end{smallmatrix}\right)$ with
$\max(0,p+q-1)\le a\le\min(p,q)$. The product coupling has $a=pq$, and the maximal coupling, of cost
$\TV(f(x),g(x))$, has $a=\min(p,q)$. Between deterministic 1-cells the coupling is unique, and its cost is
$0$ for equal 1-cells and $1$ for any two distinct ones, because distinct deterministic maps differ at some
$x$ and the cost is a maximum over $x$.
\end{proposition}

\begin{definition}[Whiskering]\label{def:whisker}
Precomposition by a deterministic $u\colon W\to B$ is $(\alpha\circ u)(w)=\alpha(u(w))$. Postcomposition by
$h\colon B'\to\D(B'')$ copies $h$ on equal pairs and runs it independently on unequal pairs:
\[
(h\cdot\alpha)(x)(z,z')=\sum_y\alpha(x)(y,y)\,h(y)(z)\,[z=z']+\sum_{y\ne y'}\alpha(x)(y,y')\,h(y)(z)\,h(y')(z').
\]
\end{definition}

This is the whiskering of \cite[Thm.~2.3(ii)]{Conditionals}. It sends diagonal couplings to diagonal
couplings. Running $h$ independently also on equal pairs would not preserve identities: for
$h=k(\tfrac9{10},\tfrac1{10})$ it sends $\tfrac12\Delta$ to
$\left(\begin{smallmatrix}0.41&0.09\\0.09&0.41\end{smallmatrix}\right)$, so it cannot be the first component
of any functor.

\begin{proposition}[Whiskering defects]\label{prop:defect}
Let $F_2(\beta,\alpha)=(h\cdot\beta)\cdot(h\cdot\alpha)-h\cdot(\beta\cdot\alpha)$.
\begin{enumerate}[label=(\roman*)]
\item Precomposition by a deterministic map is strictly functorial.
\item If $\alpha$ is an identity, then $F_2(\beta,\alpha)=0$.
\item With $f=g$ having uniform rows and $\alpha=\beta$ the swap coupling at each $x$,
  $F_2(\beta,\alpha)(x)=t\left(\begin{smallmatrix}-1&1\\1&-1\end{smallmatrix}\right)$, where $t=\tfrac{369}{2500}$ for
  $h=k(\tfrac9{10},\tfrac1{10})$ and $t=\tfrac14$ for $h$ constant and uniform.
\item Every defect has zero row and column sums. It is a tangent vector of the coupling polytope of
  $h^\#f(x)$ and $h^\#k(x)$, not a coupling.
\end{enumerate}
\end{proposition}
\begin{proof}
(i) Both sides evaluate the same couplings at $u(w)$. (ii) $h\cdot\Delta=\Delta$, and $\Delta$ is a unit for
gluing. (iii) $\beta\cdot\alpha=\mathrm{swap}\cdot\mathrm{swap}=\Delta$, so $h\cdot(\beta\cdot\alpha)$ is diagonal,
while $h\cdot\mathrm{swap}$ has mass off the diagonal; the values are exact computations. (iv) Both terms
couple the same two marginals.
\end{proof}

So the defect lives in the tangent space of a coupling polytope. There the Fisher metric, which is
twice the Hessian of Guillemin's potential \cite[Prop.~4.5]{Conditionals}, gives it a size. This is the
only sense in which the defect meets the toric geometry that we can establish.

%=====================================================================
\section{Second- and third-order states}\label{sec:levels}
%=====================================================================

Let $Y\subseteq\D B$ be a finite set of $n$ first-order states and $P,Q\in\D(Y)$ second-order states with full
support. A level-two coupling is $\Pi\in\D(Y\times Y)$ with marginals $P$ and $Q$, and $\Cpl_2(P,Q)$ is the
transportation polytope on $K_{n,n}$ with these marginals.

\begin{proposition}[Delzant at levels two and three]\label{prop:l23}
If $(P,Q)$ is nondegenerate, then $\Cpl_2(P,Q)$ is Delzant of dimension $(n-1)^2$. The same holds for
third-order states $R,S\in\D(Z)$, with $Z$ a finite set of second-order states.
\end{proposition}
\begin{proof}
\Cref{thm:delzant} uses only the marginals, not what the points of $Y$ or $Z$ are.
\end{proof}

\paragraph{Flattening.} The map
\[
\mu^*\colon\Cpl_2(P,Q)\to\Cpl(\mu_BP,\mu_BQ),\qquad
\mu^*(\Pi)(b,b')=\sum_{\rho,\rho'}\Pi(\rho,\rho')\,\rho(b)\,\rho'(b'),
\]
is affine. Its image contains $\mu_BP\otimes\mu_BQ$, the image of $P\otimes Q$.

\begin{proposition}[Properties of flattening]\label{prop:flat}
\leavevmode
\begin{enumerate}[label=(\roman*)]
\item If $Y$ consists of point masses, $\mu^*$ is the identity.
\item $\mu^*$ is not surjective in general: for $Y=\{\rho\}$ with $\rho$ uniform on two points the image is
  $\{\rho\otimes\rho\}$, while $\Cpl(\rho,\rho)$ is a segment containing the diagonal coupling.
\item For $n=2$ and $\rho_1\ne\rho_2$, $\mu^*$ is injective: the segment direction maps to
  $(\rho_1-\rho_2)\otimes(\rho_1-\rho_2)\ne0$.
\item If $(n-1)^2>(\lvert B\rvert-1)^2$, the fibres of $\mu^*$ are positive-dimensional. For example, with three
  states on two outcomes and uniform $P=Q$ the fibre has dimension $3$, so level-two couplings carry
  information that level-one couplings lose.
\item For a kernel $k\colon B\to\D(B')$, pushing $\Pi$ forward along $k^\#\times k^\#$ commutes with flattening:
  $\mu^*\big((k^\#\times k^\#)_*\Pi\big)=(k\otimes k)^\#\mu^*(\Pi)$. This is the level-two form of stochastic
  change of base \cite[Prop.~9.14]{Tower}.
\end{enumerate}
\end{proposition}
\begin{proof}
(i) $\rho(b)=[\rho=\delta_b]$. (ii) and (iii) are direct. (iv) The image lies in a space of dimension
$(\lvert B\rvert-1)^2$; the example is computed exactly. (v) Both sides equal
$\sum\Pi(\rho,\rho')\,k^\#\rho\otimes k^\#\rho'$.
\end{proof}

\paragraph{Uniform marginals.} Uniform $n\times n$ marginals are degenerate for every $n\ge2$. For $n=2$
the polytope is still a segment, so $X=\CP^1$ and it is monotone, with special point the product
coupling. For $n\ge3$ the polytope is not simple and has no smooth toric manifold.

%=====================================================================
\section{Disc potentials and optimal transport}\label{sec:potential}
%=====================================================================

For a Delzant polytope with facet functions $\ell_c$ and normals $v_c$, the potential of the torus fibre over
$u$ is $W_u(z)=\sum_cT^{\ell_c(u)}z^{v_c}$ over facets, for compact toric Fano manifolds \cite{ChoOh2006} and,
with corrections, for general compact toric manifolds \cite{FOOO2010}. A fibre with a flat connection has
nonvanishing Floer cohomology when it corresponds to a critical point. For coupling polytopes the facets
are cells, so $W$ has at most $nm$ terms, with exponents given by the normals of \cref{sec:setup}. The facets
are not the cycles of $K_{n,m}$.

\begin{proposition}[The del Pezzo potential]\label{prop:dp6}
At the uniform $2\times3$ point all facet functions equal $\tfrac16$ at the product coupling, and
$W=T^{1/6}(x+y+x^{-1}+y^{-1}+xy+x^{-1}y^{-1})$. Up to the factor $T^{1/6}$, $W$ has six critical points:
$(1,1)$, $(1,-1)$, $(-1,1)$, $(-1,-1)$, $(\zeta,\zeta)$ and $(\zeta^2,\zeta^2)$ for a primitive cube root of unity
$\zeta$. All six have $\lvert x\rvert=\lvert y\rvert=1$, so they lie over the product coupling and differ only in
the flat connection. The critical values are $6,-2,-2,-2,-3,-3$.
\end{proposition}
\begin{proof}
Exact solution of $x\partial_xW=y\partial_yW=0$; for instance $W(\zeta,\zeta)=3\zeta+3\zeta^2=-3$. The number $6$ equals twice the area $3$ of the Newton hexagon,
as Kushnirenko's theorem predicts.
\end{proof}

\begin{example}[Non-monotone cases]
\leavevmode
\begin{enumerate}[label=(\alph*),leftmargin=*]
\item \emph{A segment.} For the $2\times2$ segment $[\max(0,p+q-1),\min(p,q)]$, $X=\CP^1$ and the special fibre is
  the midpoint. That is the product coupling $a=pq$ when $p=\tfrac12$, but not in general: for
  $p=q=\tfrac13$ the midpoint is $\tfrac16$ and $pq=\tfrac19$.
\item \emph{A $3\times3$ case.} For $P=(0.40,0.35,0.25)$ and $Q=(0.30,0.43,0.27)$ the polytope is Delzant with $18$
  vertices and $9$ facets, one per cell, so $W$ has $9$ terms.
\end{enumerate}
\end{example}

\begin{proposition}[Square cases are not Fano]\label{prop:notfano}
Let $n\ge3$ and let $(f,g)$ be a nondegenerate $n\times n$ pair for which every cell is a facet. Then
$X_{f,g}$ is not Fano. In the $3\times3$ example above, $-K_X$ is nef and big, so $X$ is weak Fano.
\end{proposition}
\begin{proof}
$-K_X=\sum_cD_c$ corresponds to the polytope $\{\langle u,v_c\rangle\ge-1\ \forall c\}$, which is the transportation
polytope with all row and column sums $n$, that is $n$ times the Birkhoff polytope. If $-K_X$ were ample,
this polytope would have normal fan equal to the smooth fan of $X$ and would be simple, but the Birkhoff
polytope is not simple for $n\ge3$. For nefness, each maximal cone (a spanning tree $T$) must give a vertex
of that polytope, that is, the tree solution of $T$ with all marginals equal to $n$ must be nonnegative.
This is checked for all $18$ trees of the example. Bigness is full-dimensionality.
\end{proof}

So Cho--Oh's Fano theorem applies to the $2\times3$ family, whose surfaces are all Fano, but not to square
cases with $n\ge3$, where the general theory of \cite{FOOO2010} is needed.

\begin{proposition}[Optimal transport is a moment-map problem]\label{prop:ot}
For a cost $c\colon Y\times Y'\to\mathbb{R}$, the optimal-transport value $\min_{\alpha\in\Cpl(f,g)}\langle c,\alpha\rangle$ is
the minimum over $X_{f,g}$ of the function $\langle c,\mu\rangle$, where $\mu$ is the moment map. It is attained at a
torus-fixed point, that is, a vertex: a coupling supported on a spanning tree. If $c$ is integral on
$\Lambda$, then $\langle c,\mu\rangle$ is the Hamiltonian of a circle action.
\end{proposition}
\begin{proof}
$\mu$ maps $X_{f,g}$ onto $\Cpl(f,g)$. A linear function on a polytope is minimized at a vertex, and the
preimage of a vertex is a fixed point. A component $\langle\xi,\mu\rangle$ with $\xi$ in the lattice generates the
circle action of $\xi$ \cite{Delzant1988}. The script checks the first statement by linear programming against
vertex enumeration.
\end{proof}

This is where the polytope picture meets the obstruction theory of \cite{Tower}. Kantorovich comparisons
between second-order states \cite[Prop.~9.14]{Tower}, and distances to effective sets, are linear programs
over polytopes of this kind. Their optima are spanning-tree couplings, and their dual variables are the
certificates. These vertices are not ``committed'' states in the tower's sense. A spanning-tree coupling
has $n+m-1$ nonzero entries and is not a point mass.

%=====================================================================
\section{Stability conditions}\label{sec:stability}
%=====================================================================

A Bridgeland stability condition on a triangulated category is a pair $(Z,\mathcal A)$: the heart $\mathcal A$
of a bounded $t$-structure, and an additive central charge $Z\colon K_0\to\mathbb{C}$ that maps nonzero objects of
$\mathcal A$ into the upper half-plane together with $\mathbb{R}_{<0}$, with Harder--Narasimhan filtrations and the
support property \cite{Bridgeland2007}. The stability conditions form a complex manifold
\cite{Bridgeland2007}. On any smooth projective surface, for any ample class $\omega$, there is one with
central charge $Z_\omega(E)=-\int e^{-i\omega}\ch(E)$ and a tilted heart \cite{ArcaraBertram2013,MacriSchmidt2017}.
(Bridgeland's earlier construction \cite{Bridgeland2008} is for K3 surfaces.)

\begin{proposition}[Marginals give stability conditions]\label{prop:stab}
For $(f,g)$ in a chamber of the $2\times3$ family, $\omega_{f,g}$ defines a stability condition $\sigma_{f,g}$ on
$D^b(\Coh X_{f,g})$, depending continuously on $(f,g)$ within the chamber, with
\[
Z(\mathcal O_X)=\operatorname{area}\Cpl(f,g),\qquad Z(\mathcal O_x)=-1,\qquad
Z(\mathcal O_{D_c})=\tfrac12D_c^2+i\,\operatorname{length}\{\alpha_c=0\}.
\]
As $(f,g)$ approaches a wall from inside a chamber, one edge length tends to $0$, its divisor $E$ is a
$(-1)$-curve, and $Z(\mathcal O_E)\to-\tfrac12$.
\end{proposition}
\begin{proof}
$\omega_{f,g}$ is ample by \cref{prop:kahler2x3}. The values follow from $\ch(\mathcal O_X)=1$,
$\ch(\mathcal O_x)=[\mathrm{pt}]$ and $\ch(\mathcal O_D)=D-\tfrac12D^2$, together with \cref{prop:kahler2x3}. The
statement about walls is \cref{prop:wall} with the self-intersections of \cref{thm:chambers}, and is
checked in \cite[\S7.1]{Conditionals}.
\end{proof}

\paragraph{Mirror symmetry.} For a toric Fano manifold the disc potential is the Hori--Vafa
Landau--Ginzburg potential \cite{ChoOh2006}. The disc potential enters central charges through oscillatory
integrals $\int e^{-W/z}$ over integral cycles, whose asymptotics involve the $\widehat\Gamma$-class
\cite{Iritani2009,GGI2016}, not through its values. The critical values in \cref{prop:dp6} are the
eigenvalues of quantum multiplication by $c_1$, which is what the Gamma conjectures concern.

\begin{question}\label{q:stab}
\leavevmode
\begin{enumerate}[label=(\roman*),leftmargin=*]
\item Does the path $\sigma_{f,g}$ continue across a wall, glued along Orlov's decomposition
  $D^b(\mathrm{Bl}_pX)=\langle\mathcal O_E(-1),D^b(X)\rangle$ \cite{Orlov1992}?
\item Is there a natural object attached to an interior coupling, for instance the mirror of its torus
  fibre with a flat connection, whose central charge has a probabilistic meaning?
\item Do the quantum-corrected central charges with $\omega=\omega_{f,g}$ define stability conditions on these
  surfaces, and does the monotone point play a special role?
\end{enumerate}
\end{question}

\paragraph{Not claimed.} We make no identification of the disc potential with a central charge, of
spanning-tree couplings with a heart, of the obstruction tower with Harder--Narasimhan filtrations, or of
hallucination with instability.

%=====================================================================
\section{The $A_\infty$ question}\label{sec:ainfty}
%=====================================================================

The hom-categories of coupling 2-cells are strictly associative, so as $A_\infty$-categories they have
$m_k=0$ for $k\ge3$. If a linearization also has $m_1=0$, the arity-two $A_\infty$-functor equation reduces to
$F_1(m_2(\beta,\alpha))=m_2(F_1\beta,F_1\alpha)$, which says that whiskering is functorial. \Cref{prop:defect}(iii)
shows that it is not. Hence the defect can be the second component of an $A_\infty$-functor only in a
linearization with nonzero $m_1$, for instance chains on coupling polytopes in which the defects are
boundaries \cite[Conj.~7.3]{Conditionals}. No such linearization is known. Statements identifying the defect with
the Hessian of a disc potential, or the functor equation with a Picard--Lefschetz formula, are withdrawn
(\cref{app:changes}).

\appendix
%=====================================================================
\section{Changes from the earlier draft}\label{app:changes}
%=====================================================================

\begin{center}\small
\begin{tabular}{@{}R{0.36\textwidth}R{0.56\textwidth}@{}}
\toprule
Earlier claim & Status\\
\midrule
22 chambers; 4/12/6 quadrilaterals, pentagons, hexagons & 18 chambers; 2 triangles, 9 quadrilaterals, 6 pentagons, 1
hexagon (\cref{thm:chambers}). The accompanying script printed 18; a coordinate-order bug in its
\texttt{intersect} produced wrong vertex counts.\\\addlinespace[3pt]
Example 4.2 a quadrilateral & a triangle, $\CP^2$ (\cref{ex:2x3})\\\addlinespace[3pt]
Self-intersection table & corrected; the earlier patterns violated $\sum D_i^2=12-3k$\\\addlinespace[3pt]
Flips; type change across $f_1=\tfrac12$ & neither occurs (\cref{prop:wall})\\\addlinespace[3pt]
Costs $c(\sigma,e_0)=c(e_0,\tau)=\tfrac12$ & both are $1$; cost is a maximum over $x$\\\addlinespace[3pt]
Postcomposition running $h$ independently on equal pairs & not unital; replaced by
\cref{def:whisker}. The earlier example then has zero defect, and new examples are nonzero
(\cref{prop:defect})\\\addlinespace[3pt]
Flattening surjective; Example 10.10 constant & not surjective; for point masses it is the identity
(\cref{prop:flat})\\\addlinespace[3pt]
Uniform level two monotone iff every cell is a facet; $n=2$ not monotone & uniform $n\times n$ is degenerate; $n=2$
is monotone, $n\ge3$ is not simple\\\addlinespace[3pt]
Level-two law restricted to couplings (Thm.~10.13) & ill-typed; replaced by \cref{prop:flat}(v)\\\addlinespace[3pt]
Facets are cycles of $K_{n,n}$; 15 facets; smooth Fano fourfold & facets are cells; $9$ facets; uniform $3\times3$ is not
simple; nondegenerate square $n\ge3$ is not Fano (\cref{prop:notfano})\\\addlinespace[3pt]
Lattice distances are sums over cycles & they are the entries of the coupling at the fibre position\\\addlinespace[3pt]
Critical points of $W$ are vertices (committed states) & critical points are interior torus fibres with flat
connections (\cref{prop:dp6}); vertices are fixed points and optimal-transport optima (\cref{prop:ot})\\\addlinespace[3pt]
$F_2=\mathrm{Hess}\,W$; Picard--Lefschetz formula; ``symbolic proof'' & withdrawn (\cref{sec:ainfty}); the cited
verification compared $0.1152$ with $2$\\\addlinespace[3pt]
Cycle counts $N_{2k}=\frac1{2k}\binom nk^2k!(k-1)!$; $K_{4,4}$ total 156 & $N_{2k}=\binom nk^2k!(k-1)!/2$; $K_{4,4}$ has
$36+96+72=204$; $K_{5,5}$ has $100+600+1800+1440$\\\addlinespace[3pt]
Bridgeland 2008 for all surfaces; ``(i)--(iii) none known'' & Bridgeland 2008 is for K3; general surfaces are
\cite{ArcaraBertram2013}; (i) is answered by \cref{prop:stab}\\\addlinespace[3pt]
\bottomrule
\end{tabular}
\end{center}

%=====================================================================
\section{Computational checks}\label{app:code}
%=====================================================================

\texttt{check\_polytope\_companion.py} runs $26$ checks, all passing:
\begin{itemize}[leftmargin=*]
\item chambers, types, surfaces, the sum rule, orbits, the del Pezzo fan, wall-crossings and
  \cref{ex:2x3} (F1--F7);
\item \cref{prop:kahler,prop:kahler2x3} in all $18$ chambers (K1--K3);
\item \cref{prop:segment,prop:defect} (S1--S4);
\item \cref{prop:l23,prop:flat} (L1--L7);
\item \cref{prop:dp6}, the non-monotone examples and \cref{prop:notfano} (W1--W3);
\item \cref{prop:ot} by linear programming (T1);
\item the cycle counts of $K_{3,3}$ and $K_{4,4}$ by enumeration (C1).
\end{itemize}
Rational arithmetic is exact. The critical points are solved symbolically, and the linear programs use
floating point.

\begin{thebibliography}{99}

\bibitem{ArcaraBertram2013}
D.~Arcara and A.~Bertram.
\newblock Bridgeland-stable moduli spaces for {$K$}-trivial surfaces (with an appendix by M.~Lieblich).
\newblock \emph{Journal of the European Mathematical Society}, 15(1):1--38, 2013.

\bibitem{Conditionals}
V.~K. Awasthi.
\newblock Conditionals, couplings and toric geometry: established results and a {F}ukaya programme.
\newblock Revised draft, Johns Hopkins University, October 2026.

\bibitem{Tower}
V.~K. Awasthi.
\newblock The three-level posterior tower: hierarchical {B}ayesian inference and mixed distributive laws.
\newblock Technical report, Johns Hopkins University, October 2026. Revised version 9.

\bibitem{Bridgeland2007}
T.~Bridgeland.
\newblock Stability conditions on triangulated categories.
\newblock \emph{Annals of Mathematics}, 166(2):317--345, 2007.

\bibitem{Bridgeland2008}
T.~Bridgeland.
\newblock Stability conditions on {K}3 surfaces.
\newblock \emph{Duke Mathematical Journal}, 141(2):241--291, 2008.

\bibitem{ChoOh2006}
C.-H. Cho and Y.-G. Oh.
\newblock Floer cohomology and disc instantons of {L}agrangian torus fibers in {F}ano toric manifolds.
\newblock \emph{Asian Journal of Mathematics}, 10(4):773--814, 2006.

\bibitem{Delzant1988}
T.~Delzant.
\newblock Hamiltoniens p\'eriodiques et images convexes de l'application moment.
\newblock \emph{Bulletin de la Soci\'et\'e Math\'ematique de France}, 116(3):315--339, 1988.

\bibitem{FOOO2010}
K.~Fukaya, Y.-G. Oh, H.~Ohta, and K.~Ono.
\newblock Lagrangian {F}loer theory on compact toric manifolds {I}.
\newblock \emph{Duke Mathematical Journal}, 151(1):23--174, 2010.

\bibitem{Fulton1993}
W.~Fulton.
\newblock \emph{Introduction to Toric Varieties}.
\newblock Annals of Mathematics Studies 131. Princeton University Press, 1993.

\bibitem{GGI2016}
S.~Galkin, V.~Golyshev, and H.~Iritani.
\newblock Gamma conjectures for {F}ano manifolds.
\newblock \emph{Duke Mathematical Journal}, 165(11):2005--2077, 2016.

\bibitem{Guillemin1994}
V.~Guillemin.
\newblock Kaehler structures on toric varieties.
\newblock \emph{Journal of Differential Geometry}, 40(2):285--309, 1994.

\bibitem{Iritani2009}
H.~Iritani.
\newblock An integral structure in quantum cohomology and mirror symmetry for toric orbifolds.
\newblock \emph{Advances in Mathematics}, 222(3):1016--1079, 2009.

\bibitem{MacriSchmidt2017}
E.~Macr\`i and B.~Schmidt.
\newblock Lectures on {B}ridgeland stability.
\newblock In \emph{Moduli of Curves}, Lecture Notes of the Unione Matematica Italiana 21. Springer,
2017. arXiv:1607.01262.

\bibitem{Orlov1992}
D.~O. Orlov.
\newblock Projective bundles, monoidal transformations, and derived categories of coherent sheaves.
\newblock \emph{Izvestiya Rossiiskoi Akademii Nauk. Seriya Matematicheskaya}, 56(4):852--862, 1992.

\end{thebibliography}
\end{document}
